
El \cref{chap:cap3} estableció la formulación mecanicista del metabolismo larval mediante el marco de Presupuestos Energéticos Dinámicos (DEB). Como se demostró, el sistema de ecuaciones diferenciales ordinarias resultante —que acopla el sustrato $S(t)$, la biomasa estructural $B(t)$, los lípidos de reserva $L(t)$ y la producción de \ce{CO2} $C(t)$— captura la dinámica bioenergética del proceso de bioconversión con \textit{Hermetia illucens}. Sin embargo, la implementación operativa de dicho modelo revela una limitación estructural crítica: la variable de estado central, $B(t)$, no es observable en tiempo real. En la práctica experimental, la biomasa larval solo puede cuantificarse mediante muestreo destructivo con intervalos típicos de \num{5} días, lo que convierte al modelo DEB en un predictor de lazo abierto durante la mayor parte del ciclo productivo.

Esta condición de \emph{estado parcialmente observable} constituye el punto de partida del presente capítulo. Cuando el modelo DEB opera sin retroalimentación dinámica de $B(t)$, los errores de predicción de $\dot{m}_{\mathrm{CO_2}}^{\mathrm{DEB}}(t)$ se acumulan de manera monótona, particularmente ante perturbaciones no modeladas como cambios en la composición del sustrato, fluctuaciones térmicas o variaciones en la densidad larval efectiva. La literatura reciente \parencite{kobelskiModelbasedProcessOptimization2024, padmanabhaModellingOptimalControl2023} ha documentado que, en sistemas biotecnológicos de insectos, la desconexión entre el estado biológico real y su estimación teórica es una de las principales fuentes de incertidumbre predictiva.

Para cerrar esta brecha, se propone una arquitectura híbrida que integra el modelo mecanicista DEB con un estimador externo de biomasa basado en Procesos Gaussianos (GP). La lógica subyacente es la siguiente: si bien $B(t)$ no puede medirse directamente con alta frecuencia temporal, las variables operativas del reactor —concentración de \ce{CO2} $C(t)$ en ppm, temperatura $T$ en °C, humedad relativa HR en \% y el tiempo transcurrido desde el inicio del ciclo $t$ en días— sí se registran de manera continua mediante sensores NDIR de bajo costo. Estas variables constituyen un \emph{proxy} indirecto del estado metabólico del sistema, dado que la respiración aeróbica, la temperatura del sustrato exotérmico y la humedad intersticial están todas moduladas, en última instancia, por la biomasa activa presente en el reactor \parencite{rossiEstimatingDynamicsGreenhouse2024, biagiDevelopmentMachineLearningbased2024}.

La elección del Proceso Gaussiano como estimador no es arbitraria. En contextos con datos experimentales limitados —característicos de ciclos de bioconversión de \num{14} a \num{20} días donde el número de réplicas es económicamente restrictivo— los GP ofrecen ventajas distintivas sobre arquitecturas de aprendizaje profundo: (i) proporcionan una distribución predictiva completa, no solo un punto estimado, lo que permite propagar la incertidumbre de $\hat{B}(t)$ hacia la predicción de emisiones; (ii) su base en kernels de covarianza permite capturar no linealidades complejas sin requerir grandes volúmenes de entrenamiento; (iii) la regularización inherente al margen de ruido $\sigma_n^2$ mitiga el riesgo de sobreajuste, crítico cuando el número de mediciones destructivas de biomasa es reducido \parencite{raissiPhysicsinformedNeuralNetworks2019, adadiPeekingBlackBoxSurvey2018}. Las condiciones de convergencia del estimador GP en este contexto —continuidad de la función media, definida positividad del kernel y acotamiento de la varianza de ruido— se satisfacen bajo los supuestos operativos del reactor; los detalles formales se remiten al Anexo~C.

Desde la perspectiva del modelado de caja gris (\emph{grey-box modelling}), la arquitectura DEB+GP responde a una jerarquía epistémica deliberada: el modelo mecanicista DEB aporta la \emph{tendencia base} interpretable, derivada de principios bioenergéticos de conservación de masa, mientras que el GP corrige el \emph{error residual} asociado a la variabilidad estocástica del sistema y a la imposibilidad de medir $B(t)$ en continuo. Esta división de labor preserva la transparencia física del modelo —cada parámetro DEB conserva su significado biológico explícito— al tiempo que otorga al sistema la adaptabilidad empírica necesaria para operar en condiciones agroindustriales reales \parencite{weichertReviewMachineLearning2019}.

El flujo de información de la arquitectura propuesta opera en cascada. Primero, el estimador GP infiere la biomasa húmeda total del reactor $\hat{B}(t)$ a partir del vector de entradas operativas $\mathbf{x}(t) = [C(t), T, \mathrm{HR}, t]^\intercal$. Segundo, este pseudo-observable alimenta el modelo DEB, que calcula la tasa teórica de emisión $\dot{m}_{\mathrm{CO_2}}^{\mathrm{DEB}}(t)$. Tercero, la predicción mecanicista se confronta con la pendiente observada del sensor $\Delta C/\Delta t$ en ventanas cerradas de \num{30} min, generando un residuo operativo que guía la calibración dinámica de parámetros libres. En este esquema, el \ce{CH4} no participa en las ecuaciones de estado —pues el modelo DEB asume aerobiosis estricta— sino que opera como clasificador binario de anaerobiosis transitoria, activando alertas cuando su concentración supera el umbral de \num{100} ppm.

Las métricas de éxito de esta integración se definen de manera cuantitativa. El sistema híbrido se considerará validado si alcanza un coeficiente de determinación $R^2 \geq \num{0.70}$ en la predicción de la pendiente de emisión de \ce{CO2}, un error porcentual absoluto medio inferior al del modelo lineal estático de referencia, y una sensibilidad (Recall) superior al \num{80}\% en la detección de eventos de \ce{CH4} asociados a anaerobiosis. Estos umbrales, alineados con los criterios del \cref{chap:marco_metodologico}, garantizan que la herramienta no solo predice, sino que proporciona evidencia verificable para la gestión operativa.

El resto de este capítulo desarrolla rigurosamente cada componente de la arquitectura híbrida. La \cref{sec:cap4_arquitectura} describe el flujo de información y las interfaces entre módulos. La \cref{sec:cap4_estimador_gp} formaliza el estimador de biomasa mediante GP, incluyendo la justificación de su ubicación fuera de la ODE, la codificación del estado del reactor y la estrategia de validación \emph{leave-one-replica-out}. La \cref{sec:cap4_calibracion} detalla la calibración del modelo DEB en dos fases: semillas de parámetros a partir de \textcite{eriksenDynamicModellingFeed2022} y ajuste local mediante optimización restringida. La \cref{sec:cap4_perdida} define la función de pérdida que confronta pendientes modelo-sensor y el tratamiento residual del \ce{CH4}. La \cref{sec:cap4_validacion} establece las métricas predictivas y el análisis de residuos. Finalmente, la \cref{sec:cap4_resultados} presenta los resultados cuantitativos de la integración y la \cref{sec:cap4_conclusiones} sintetiza las implicaciones para el sistema de monitoreo.
